\section{Evaluation}
\label{sec:evaluation}

The evaluation measures five properties: bypass coverage, corrective-feedback
effectiveness, false-positive rate, runtime overhead, and coverage over
real-world contracts.

\subsection{Experimental Setup}

% TODO: Fill in agent versions, kernel version, hardware, policy set.

\subsection{Bypass Coverage (RQ1)}

Does OS-level enforcement detect the same contract violation regardless of
the tool path through which the agent triggers it?

We execute the same forbidden operation through five paths: (1)~direct tool
call, (2)~nested shell string (\texttt{bash~-c}), (3)~Python
\texttt{subprocess}, (4)~compiled helper binary, and (5)~direct
\texttt{execve} from generated C code. For each path, we record whether
ActPlane detects the violation and whether the corrective feedback payload
is delivered to the agent.

% TODO: Insert coverage matrix table.

\subsection{Feedback and Recovery (RQ2)}

Does corrective feedback improve agent task completion and reduce repeated
violations compared with enforcement without feedback?

We run agent tasks under four conditions: (C1)~prompt-only, in which the
contract is stated in the prompt but not enforced; (C2)~audit, in which ActPlane
reports violations but does not fail operations; (C3)~enforcement without
feedback, in which the operation fails with \texttt{-EPERM} or \texttt{SIGKILL}
but no structured reason; (C4)~enforcement with feedback, in which the operation
fails and ActPlane injects the rule name, reason, and remediation into the
agent's context.

The primary comparison is C3 vs.\ C4: the marginal benefit of corrective
feedback. The C1-to-C2 comparison measures how often prompt-only contracts
are violated. The C2-to-C3 comparison isolates the effect of enforcement
from observation.

% TODO: Insert results table.

\subsection{False Positives (RQ3)}

Do benign tasks and sanctioned contract paths (declassification, gate
mediation, temporal gates) proceed without spurious enforcement?

% TODO: Insert false-positive results.

\subsection{Overhead (RQ4)}

What is the per-event and end-to-end cost of label propagation and rule
checking?

We measure per-event latency for fork, exec, open, write, and connect
operations with and without ActPlane attached, and report p50 and p99
latencies. We also measure wall-clock time for representative agent tasks
and BPF map memory consumption as a function of rule count and active
process count.

% TODO: Insert latency table and overhead figures.

\subsection{Corpus Coverage (RQ5)}

How many real-world contracts from the corpus study
(\S\ref{sec:corpus}) are enforceable by ActPlane?

We apply the D1--D7 coding dimensions to each extracted contract and
report the fraction that is observable at the syscall boundary (D3),
expressible in the ActPlane DSL (D4), enforceable without false positives
(D5), and temptable in a realistic agent task (D7). We report inter-rater
agreement (Cohen's $\kappa$) for each dimension.

% TODO: Insert corpus coverage funnel table.


\section{Related Work}
\label{sec:related}

\subsection{In-Kernel Provenance and Information Flow}

PASS introduced provenance as a storage-layer
primitive~\cite{muniswamy2006pass}. Hi-Fi and LPM used LSM hooks to capture
trustworthy whole-system provenance~\cite{pohly2012hifi,bates2015lpm}.
CamFlow extended this to a streaming provenance framework over processes,
files, and sockets~\cite{pasquier2017camflow}.

CamQuery is the closest prior
system~\cite{pasquier2018camquery}. It executes provenance queries inline
in the kernel, propagates labels across process, file, and network edges,
and can deny operations that violate a policy. ActPlane shares this
label-propagation model but differs in three respects: it targets
cooperative AI agents rather than adversarial intrusions, compiles an
agent-oriented source/sink DSL rather than general provenance queries, and
returns corrective feedback to the violating actor. The implementation
differs as well: ActPlane uses the eBPF/BPF-LSM substrate rather than a
loadable kernel module.

TaintDroid, Dytan, libdft, and Panorama established dynamic taint tracking
at runtime, binary-instrumentation, and emulation
layers~\cite{enck2010taintdroid,clause2007dytan,kemerlis2012libdft,yin2007panorama}.
ActPlane operates at object granularity rather than byte granularity,
trading precision for low-overhead whole-system coverage.

\subsection{eBPF and LSM Enforcement}

BPF-LSM, originally proposed as KRSI, allows eBPF programs to attach to
Linux security hooks and return allow/deny
verdicts~\cite{singh2019krsi,linuxbpflsm}. ActPlane uses BPF-LSM as its
pre-operation enforcement backend. Contemporary eBPF enforcement tools
such as Tetragon, Tracee, and eBPF-PATROL provide per-event predicate
matching and, in Tetragon's case, a boolean lineage flag
(\texttt{followChildren}) that propagates across fork/exec. These systems
evaluate per-event predicates or single-channel lineage flags rather than
cross-channel labeled information
flow~\cite{ciliumtetragon,aquasecuritytracee,ghimire2025ebpfpatrol}.

\subsection{Agent Guardrails and Information-Flow Control}

AgentSpec and Progent provide deterministic runtime enforcement over
tool-call names, arguments, and framework
actions~\cite{wang2025agentspec,shi2025progent}. These systems enforce at
the tool-call boundary; ActPlane complements them by enforcing below the
tool API, where effects that bypass the framework are still visible.

FIDES and CaMeL apply typed information-flow control and capability
separation within the agent loop or
scaffolding~\cite{costa2025fides,debenedetti2025camel}. ActPlane applies
the same style of information-flow reasoning at the OS boundary, where
enforcement persists across context windows and cannot be circumvented by
subprocess indirection.

\subsection{Agent Instruction Files}

Recent empirical studies characterize \texttt{CLAUDE.md},
\texttt{AGENTS.md}, and related agent instruction files along dimensions
of content taxonomy, maintenance practices, and effect on task
efficiency~\cite{chatlatanagulchai2025manifests,chatlatanagulchai2025agentreadmes,
santos2025claudeconfig,lulla2026agentsmd}. ActPlane's corpus study
(\S\ref{sec:corpus}) asks a different question: which instructions
constitute cross-object information-flow contracts, and which of those can
be enforced at the OS level?


\section{Conclusion}
\label{sec:conclusion}

ActPlane provides a programmable OS-level control plane for agent harnesses.
It compiles agent-declared behavioral contracts into eBPF/BPF-LSM labeled
information-flow rules, propagates labels across process, file, and endpoint
objects at kernel hooks, and returns kernel-detected violations as semantic
feedback to the agent. By supporting audit, pre-operation denial, and
harness-level termination with structured remediation, ActPlane enables
agents to recover from contract violations rather than fail opaquely.
ActPlane unifies labeled information-flow observation and enforcement with
an agent-oriented contract language and a feedback loop on the eBPF/BPF-LSM
substrate, bridging the intent--behavior gap that single-level harnesses
leave open.
